\section{Formalizing Check-Effect Bindings}
\label{sec:formalization}

In this section, we present the formal semantics of the Agent Workflow Intermediate Representation with Bindings (\textsf{AWIR-B}), define precondition assertions and binding mechanisms, and state our relative soundness theorem.

\subsection{AWIR-B Syntax and Representation}
An agent workflow is modeled as a directed acyclic graph $W = (V, E, C)$, where:
\begin{itemize}
    \item $V = \{n_1, n_2, \dots, n_k\}$ is a finite set of workflow nodes, representing tool calls, decision gates, or state transformations.
    \item $E \subseteq V \times V$ is a set of directed control-flow and data-dependency edges.
    \item $C: V \to \mathcal{T}$ maps each node $n_i$ to a tool contract $T \in \mathcal{T}$.
\end{itemize}

\begin{definition}[Tool Contract]
A tool contract $T$ is a tuple $T = (\text{op}, \text{res}, \text{pred}, \text{mech}, \text{auth})$, where:
\begin{itemize}
    \item $\text{op} \in \{\text{READ}, \text{CHECK}, \text{EFFECT}, \text{WRITE}\}$ specifies the semantic operation type.
    \item $\text{res} \in \mathcal{R}$ identifies the target resource argument.
    \item $\text{pred}$ is a boolean predicate asserting a condition over $\text{res}$ (valid when $\text{op} = \text{CHECK}$).
    \item $\text{mech} \in \{\text{NONE}, \text{ATOMIC}, \text{VERSION}, \text{ETAG}, \text{STABLE\_IDENTITY}, \text{RECHECK}, \text{LOCK}\}$ defines the concurrency protection mechanism.
    \item $\text{auth} \in \{\text{True}, \text{False}\}$ indicates whether the binding mechanism is enforced by the authoritative resource server.
\end{itemize}
\end{definition}

\subsection{Precondition Assertions \& Protection Invariants}
When a node $n_c$ with operation $\text{op}(n_c) = \text{CHECK}$ executes on resource $r$, it introduces a \emph{propagated resource assertion} $\alpha = (r, \phi, v_c, n_c)$, where $\phi$ is the predicate evaluated at version $v_c$ of resource $r$.

A downstream node $n_e$ with operation $\text{op}(n_e) = \text{EFFECT}$ operating on $r$ depends on assertion $\alpha$ if there exists a directed path from $n_c$ to $n_e$ in $W$.

\begin{definition}[Authority-Enforced Binding]
A binding mechanism $\text{mech}$ protecting a dependency relationship $(n_c, n_e)$ on resource $r$ is \emph{Authority-Enforced} if and only if:
\begin{equation}
    \text{IsProtected}(n_c, n_e) \iff \text{auth}(n_e) = \text{True} \land \text{mech}(n_e) \neq \text{NONE}
\end{equation}
Mechanisms such as \textsf{ATOMIC} transactions, matching HTTP \textsf{ETAG} validation, or server-side conditional \textsf{VERSION} checks satisfy authority enforcement because the underlying resource authority atomically rejects the effect if $r$ was mutated post-$n_c$.
\end{definition}

\subsection{Definition of Check-Effect Binding Vulnerability}
\begin{definition}[Missing Check-Effect Binding]
Let $W = (V, E, C)$ be a workflow graph. A \emph{Missing Check-Effect Binding} vulnerability exists at path $(n_c \rightsquigarrow n_e)$ if:
\begin{enumerate}
    \item $\text{op}(n_c) = \text{CHECK}$ and $\text{op}(n_e) = \text{EFFECT}$.
    \item Both nodes reference the same target resource: $\text{res}(n_c) = \text{res}(n_e) = r$.
    \item There exists a valid reachability path $n_c \rightsquigarrow n_e$ in $W$.
    \item The propagated binding set along $n_c \rightsquigarrow n_e$ contains no authority-enforced binding:
    \begin{equation}
        \forall b \in \text{Bindings}(n_c \rightsquigarrow n_e), \quad \text{auth}(b) = \text{False} \lor \text{mech}(b) = \text{NONE}
    \end{equation}
\end{enumerate}
\end{definition}

\subsection{Relative Soundness Proof Sketch}
\begin{theorem}[Relative Soundness of Topological Dependency Propagation]
\label{thm:soundness}
Let $W = (V, E, C)$ be an acyclic workflow graph with complete and accurate tool contracts $\mathcal{T}$. If AgentLint reports zero findings for $W$, then for every execution trace $\sigma$ of $W$ under any arbitrary external state mutation schedule $S$, no effect node $n_e$ can execute against an invalidated check predicate $\phi(r)$.
\end{theorem}

\begin{proof}[Proof Sketch]
We proceed by structural induction on the topological ordering of nodes in $W$, established via Kahn's algorithm~\cite{kahn1962topological}.
\paragraph{Base Case} For the initial node $n_0$, the propagated resource and binding sets are empty, and no effect nodes exist.
\paragraph{Inductive Step} Assume for all predecessor nodes $p \in \text{Pred}(n_k)$, the propagated resource set $\mathcal{P}_{\text{res}}(p)$ and binding set $\mathcal{P}_{\text{bind}}(p)$ correctly represent all reachable check assertions and active bindings. At node $n_k$:
1. If $\text{op}(n_k) = \text{EFFECT}$ on resource $r$, AgentLint inspects the incoming joined sets $\bigcup_{p} \mathcal{P}_{\text{res}}(p)$ and $\bigcup_{p} \mathcal{P}_{\text{bind}}(p)$.
2. If an assertion $\alpha = (r, \phi, v_c, n_c)$ is present without a matching authority-enforced binding in $\mathcal{P}_{\text{bind}}$, a finding is generated.
3. Therefore, if zero findings are returned, every reachable check assertion on $r$ is matched by an authority-enforced binding, which guarantees server-side atomicity or rejection under schedule mutation $S$.
\end{proof}

% TODO: [Formal Methods] Complete full mechanized proof in Coq/Lean for loop-unrolled AWIR-B operational semantics.
